\chapter{Energiekalibration der Detektoren}

\section{Durchführung der Messungen}

Es wurden nun jeweils mit dem NaI-Detektor und dem HPGe-Detektor Messungen der drei Proben sowie jeweils ein Untergrundspektrum für den Detektor ohne Probe aufgenommen. Als Messdauer wurde für alle Messungen mit dem HPGe-Detektor $t_{\text{HPGe}}=\SI{300(0.5)}{\second}$ gewählt, für alle Messungen mit dem NaI-Detektor $t_\text{NaI}=\SI{400(0.5)}{\second}$. Als Verstärkungsfaktor wurde wie in Abschnitt \ref{sec:verstaerkung} erläutert, $1,700$ gewählt und konstant gehalten. Bei den Messungen wurde darauf geachtet, dass die angezeigte Zählrate vom MC2-Analyser-Programm, über welches die Daten des Detektors aufgenommen wurden, nicht deutlich über $\SI{2}{\kilo\hertz}$ betrug, um die Energiekalibration nicht zu verschlechtern \cite[4]{521versuchsanleitung} (linearer Zusammenhang von Impulshöhe zu Kanalnummer wäre dann nicht mehr gegeben\cite[291]{leotechniques}).
Die aufgenommenen Spektren sind in Abbildung \ref{fig:undergrd_nai_with_eb}-\ref{fig:eu_ge_with_eb} dargestellt. Hierbei wurde der statistische Fehler der gemessenen Anzahl $N$ (gemäß dem Poisson-Statistik-Verhalten des Fehlers von Zählexperimenten wie den durchgeführten) als $\Delta N=\sqrt{N}$ angenommen. Aus Lesbarkeitsgründen wurden diese Fehler in allen folgenden graphischen Auftragungen der Spektren weggelassen, aber in den Berechnungen gemäß Gaußscher Fehlerfortpflanzung \cite[8-9]{gerthsen} weiter beachtet. Da keine Bleiabschirmung gegen Hintergrundsignale verwendet wurde, um Rückstreustrahlung an der Bleiumgebung möglichst zu vermeiden, muss das jeweilige gemessene Untergrundspektrum (gleich für alle Probemessungen mit einem Detektor, da die Messungen hinreichend zeitnah zueinander durchgeführt wurden) von den Probenspektren abgezogen werden. Eine zeitliche Normierung muss nicht erfolgen, da gleiche Messdauern gewählt wurden. Es ist zu beachten, dass in den Histogrammen der Messdaten vom MCA-Programm jeweils der erste und letzte Kanal als \textit{overflow}-Kanal für Energien unterhalb und überhalb des vom Kanalbereich mit gewählter Verstärkung erfassten Energiebereichs genutzt werden und daher nicht eine sinnvolle Anzahl an Ereignissen mit ihrem eigentlichen eigenem Wert anzeigen. Daher werden sie bei den Analysen im folgenden immer nicht beachtet.

Die gewählten Abstände der Quellen (gewählt um eine Zählrate im gewünschten Bereich zu erhalten) sind in Tabelle \ref{tab:distance} aufgelistet.

\begin{table}[h]
    \centering
    \begin{tabular}{|c||c|c|c|c|} \hline
    Detektor & Abstand / m & & & \\ \hline
	& für & Cs & Co & Eu \\
	\hline \hline
	HPGe & & $\num{0.0850(50)}$ & $\num{0.0020(50)}$ & $\num{0.1400(50)}$ \\
	NaI & & $\num{0.0900(50)}$ & $\num{0.0100(50)}$ & $\num{0.1700(50)}$ \\
	\hline
    \end{tabular}
    \caption{Abstände der Proben bei den Messungen}
    \label{tab:distance}
\end{table}

Die Spektren entsprechen in ihrer Form alle den Erwartungen an ein Gammaspektrum. Sie enthalten eine variable Anzahl an deutlichen Totalenergie-Peaks (Photopeaks) sowie die unregelmäßige Form des Compton-Untergrunds pro Photopeak durch Compton-Streuung der $\gamma$-Photonen in variablen Streuwinkeln an Detektormaterial-Hüllenelektronen sowie den deutlich erkennbaren Abfall der Compton-Kante (bei dem maximalen Compton-Streuwinkel von $\theta=\SI{180}{\degree}$, bei dem gerade zurückgestreut wird) \cite[526-531]{gammahandbook}. Es ist aufgrund des betrachteten Energiebereichs bis ca. $\SI{2000}{\kilo\eV}$ nicht mit Paarbildung zu rechnen, es sind in den Spektren (auch bei weiterer Analyse wird sich dies bestätigen) keine deutlichen Escape-Peaks zu sehen.


\section{Energiekalibration des NaI-Detektors}

Es wird davon ausgegangen, dass die Kanäle des NaI-Detektors für die beobachteten Zählraten linear mit der Energie der $\gamma$-Photonen zusammenhängen \cite[5]{521versuchsanleitung}. Um diese Proportionalität zu überprüfen und den spezifischen funktionalen Zusammenhang für den verwendeten Detektor zu bestimmen, wurden in den gemessenen Spektren von $\ce{^60Co}, \ce{^137Cs}$ und $\ce{^152Eu}$ mit der bereits beschriebenen Methode die Gaußpeaks identifiziert. Beim Spektrum von $\ce{^152Eu}$ wurden nur die deutlichsten, klar getrennten Linien in die Kalibration einbezogen.

Auffallend ist hier zunächst im NaI-Spektrum von $\ce{^137Cs}$ eine zweite Linie, die so aufgrund des Zerfallsschemas in \ref{fig:zerfallsschema_cs} nicht erwartet wurde, da die zweite gelistete Übergangsenergie von $E_\gamma=\SI{283.5(1)}{\kilo\eV}$ sehr unwahrscheinlich ist. Vermutlich handelt es sich hier um eine Röntgenlinie, bei der keine Kernanregung zu einer Photonemission geführt haben, sondern eine Anregung eines Hüllenelektrons von $\ce{^137Cs}$. Diese Linie wurde bei der Energiekalibration nicht beachtet, da sie nicht zuverlässig einer Energie zugeordnet werden konnte, bei Eintragung auf der finalen Anpassungsfunktion als die zweite $\gamma$-Photon-Übergangslinie passte sie auch deutlich sichtbar nicht zum restlichen Verlauf und wurde daher nicht weiter einbezogen.

Die Schwerpunkte der Linien der verschiedenen Nuklide und die jeweiligen Energiezuordungen nach den Angaben in \cite[8]{521versuchsanleitung} sind in Tabelle \ref{tab:zuordnungen_nai} gelistet.

\todo{Tabelle Zuordnungen!!!}

\begin{table}[h]
    \centering
    \begin{tabular}{|c|c|c|c|c|} \hline
    Nuklid & $\mu$ / Kanal & zugeordnete Energie / keV \\
	\hline
	\hline
    \end{tabular}
    \caption{Bestimmte Peakschwerpunkte $\mu$ und zugeordnete Linien der verschiedenen Nuklide für den NaI-Detektor.}
    \label{tab:zuordnungen_nai}
\end{table}

Die zugeordneten Energien werden nun gegen die bestimmten Kanäle der Peakschwerpunkte der Messungen aufgetragen. Der Zusammenhang wird als linear erwartet (s.o.) und eine Anpassungsgerade der Form
\begin{equation}
E(\text{Kanal})=a\cdot \text{Kanal}+b
\end{equation}
wird an die Datenpunkte angepasst. Für die Anpassung wird die Methode der kleinsten gemeinsamen Quadrate mithilfe der Funktion \texttt{scipy.optimize.curve\_fit} in \texttt{python} genutzt. Die Anpassungsfunktion und die eingetragenen Datenpunkte sind in Abbildung \ref{fig:nai_calibration} dargestellt. Der Wert des Bestimmtheitsmaßes $R^2=0,999$ zeigt die sehr gute Übereinstimmung der Daten mit der linearen Anpassungspunktion, der erwartete lineare Zusammenhang kann somit bestätigt werden. Als Anpassungparameter ergaben sich:

$a = \SI{120.3(1.3)}{\eV}$ und $b = \SI{0.0(9.7)e+03}{\eV}$

\jafps{nai_calibration}{Für NaI-Detektor gemessene Linien mit zugeordneten Literaturwerten für Energien sowie lineare Anpassungsgerade mit $a = \SI{120.3(1.3)}{\eV}, b = \SI{0.0(9.7)e+03}{\eV}$.}{0.6}

Damit ist die Energiekalibrationsfunktion für den NaI-Detektor (zur Abkürzung mit $x$ für den Kanal):
\begin{equation}\label{eq:kalibration_nai}
    E(x)=(\SI{120.3(1.3)}{\eV})x +\SI{0.0(9.7)e+03}{\eV}
\end{equation}
Somit können Peakschwerpunkte jetzt in Energien umgerechnet werden und in $\gamma$-Spektren von Proben unbekannter Zusammensetzungen Energien von beobachteten Linien und so die Zusammensetzung bestimmt werden.

\section{Energiekalibration des HPGe-Detektors}

Auch beim HPGe-Dektektor soll nun die lineare Proportionalität zwischen den Kanälen und der korrespondieren Impulshöhe (und damit Energie des $\gamma$-Photons, welches gemessen wird) untersucht werden um eine Kalibrationsfunktion zu erhalten. Wie für den NaI-Detektor werden aus Literaturangaben (Energien aus \cite[8]{521versuchsanleitung}) den Peakschwerpunkten der gut erkennbaren Linien der drei mit dem HPGe-Detektor gemessenen Spektren (mit Untergrundspektrum abgezogen) Energien von $\gamma$-Zerfällen zugeordnet. Die Zuordnungen sind in Tabelle \ref{tab:zuordnungen_ge} dargestellt. Aus dem Spektrum von $\ce{^152Eu}$ werden auch hier nur die deutlich erkennbaren und klar zuordbaren Linien betrachtet.

\begin{table}[h]
    \centering
    \begin{tabular}{|c|c|c|c|c|} \hline
    Nuklid & $\mu$ / Kanal & zugeordnete Energie / keV \\
	\hline
	\hline
    \end{tabular}
    \caption{Bestimmte Peakschwerpunkte $\mu$ und zugeordnete Linien der verschiedenen Nuklide für den HPGe-Detektor.}
    \label{tab:zuordnungen_ge}
\end{table}
